%  LaTeX support: latex@mdpi.com
%=================================================================
\documentclass[biology,article,submit,pdftex,moreauthors]{Definitions/mdpi}
%=================================================================
% MDPI internal commands
\firstpage{1}
\makeatletter
\setcounter{page}{\@firstpage}
\makeatother
\pubvolume{1}
\issuenum{1}
\articlenumber{0}
\pubyear{2026}
\copyrightyear{2026}
\datereceived{ }
\daterevised{ }
\dateaccepted{ }
\datepublished{ }

%=================================================================
\usepackage{longtable}
\usepackage{placeins}
\raggedbottom
\setlength{\headheight}{20pt}

% Tighter float spacing
\setlength{\textfloatsep}{8pt plus 2pt minus 2pt}
\setlength{\floatsep}{8pt plus 2pt minus 2pt}
\setlength{\intextsep}{8pt plus 2pt minus 2pt}
\setlength{\abovecaptionskip}{4pt}
\setlength{\belowcaptionskip}{2pt}

\AtBeginDocument{%
  \setlength{\abovedisplayskip}{6pt plus 2pt minus 2pt}%
  \setlength{\belowdisplayskip}{6pt plus 2pt minus 2pt}%
  \setlength{\abovedisplayshortskip}{3pt plus 1pt minus 1pt}%
  \setlength{\belowdisplayshortskip}{3pt plus 1pt minus 1pt}%
}

\titlespacing{\section}{0pt}{8pt}{2pt}
\titlespacing{\subsection}{0pt}{8pt}{2pt}
\titlespacing{\subsubsection}{0pt}{8pt}{2pt}
\titlespacing{\paragraph}{0pt}{6pt}{2pt}

%=================================================================
% Figure paths: benchmark subplots (Figs 1,3-9) and experiment results (Fig 10)
\graphicspath{{../../benchmarks/paper_figures/topic/}{../../experiments/results/topic/}}

%=================================================================
\Title{Flow-Matching-Refined Dirichlet-Prior Autoencoders for Interpretable and Structurally Balanced Single-Cell Representation Learning}

\Author{Zeyu Fu $^{1,*}$}

\AuthorNames{Zeyu Fu}

\address{%
$^{1}$ \quad State Key Laboratory of Trauma and Chemical Poisoning, Institute of Combined Injury, Chongqing Engineering Research Center for Nanomedicine, College of Preventive Medicine, Army Medical University, Chongqing, China}

\corres{Correspondence: fuzeyu09@gmail.com}

%=================================================================
\abstract{We introduce flow-matching-refined Dirichlet-prior variational autoencoders (Topic-FM) for single-cell transcriptomics. A conditional optimal-transport flow field is trained in pre-softmax $\mathbb{R}^K$ space to sharpen the logistic-normal posterior without breaking the simplex constraint, yielding latent representations that simultaneously improve label concordance, cluster geometry, and biological interpretability. Four architectural variants---Topic-FM-Base (MLP), Topic-FM-Transformer (self-attention), Topic-FM-Contrastive (MoCo), and Topic-FM-GAT (graph attention over $k$NN)---are evaluated on 56 scRNA-seq datasets. Topic-FM-Transformer attains the best composite score (NMI\,+\,ARI\,+\,ASW\,/\,3) of 0.502, a 15.6\% improvement over the prior-free Pure-VAE baseline (0.434), with concordant gains in NMI (+8.2\%), ARI (+20.4\%), and ASW (+21.7\%). Wilcoxon signed-rank tests confirm statistical significance with medium-to-large Cliff's $\delta$ effects. In a full external comparison against 23 baselines, Topic-FM-Contrastive wins 86.4\% of core-metric comparisons ($p{<}0.05$). Biological validation confirms that the learned topics correspond to coherent GO programs through both perturbation-based and decoder-$\beta$-based enrichment pathways. The flow matching refinement, combined with the Dirichlet prior's built-in interpretability, establishes Topic-FM as a practical default for single-cell representation learning.}

\keyword{single-cell RNA-seq; variational autoencoder; Dirichlet prior; flow matching; optimal transport; topic model; representation learning; interpretability}

%=================================================================
\begin{document}

\section{Introduction}\label{sec:intro}

Single-cell RNA sequencing demands representation learning methods that preserve biological structure while supporting downstream interpretation of learned features. Standard variational autoencoders (VAEs) impose isotropic Gaussian priors that place no structure on the latent space and provide no direct route from latent dimensions to gene programs.

A Dirichlet-prior (topic) VAE replaces the Gaussian prior with a logistic-normal approximation to a Dirichlet distribution, constraining latent vectors to a probability simplex ($\Delta^{K-1}$). This design yields three concrete advantages. \emph{First}, each latent dimension represents a soft topic allocation interpretable as a cell's membership weight across $K$ gene programs. \emph{Second}, the decoder contains an explicit topic--word matrix $\beta$ that directly maps topics to gene-expression probabilities, enabling biological annotation without post-hoc importance scoring. \emph{Third}, the Dirichlet prior's sparsity-promoting structure encourages compact, well-separated clusters that simultaneously improve label concordance and cluster geometry relative to prior-free baselines.

While the Dirichlet prior provides strong structural inductive bias, the logistic-normal variational posterior may under-represent fine-grained geometric structure in the pre-softmax space. We address this gap with \emph{conditional flow matching} (CFM): a lightweight velocity field trained via optimal-transport interpolation in $\mathbb{R}^K$ refines the posterior samples at inference time, sharpening cluster boundaries without breaking the simplex constraint. This flow matching module activates after a warmup period and adds negligible computational overhead.

We evaluate four flow-matching-refined Dirichlet VAE variants---Topic-FM-Base (feedforward), Topic-FM-Transformer (self-attention), Topic-FM-Contrastive (MoCo contrastive), and Topic-FM-GAT (graph attention)---across a refreshed 56-dataset catalogue. Topic-FM-Transformer achieves the highest composite score (0.502, +15.6\% over Pure-VAE), with concordant gains in NMI, ARI, and ASW. This balanced improvement, combined with built-in interpretability via $\beta$ weights and the best downstream $k$NN classification accuracy, positions the flow-matching-refined Dirichlet prior as a practical default for single-cell representation learning.

\section{Methods}\label{sec:methods}

\subsection{Model Architecture}
We evaluate four flow-matching-refined Dirichlet VAE variants, each paired with a matched prior-free baseline:

\textbf{Topic-FM-Base} (feedforward encoder, 2-layer MLP $\to$ 128 $\to$ 128) maps input counts to a logistic-normal distribution over $K{=}10$ topics: $\log\theta = \mu + \sigma \cdot \epsilon$, $\theta = \text{softmax}(\log\theta)$. A learned decoder matrix $\beta \in \mathbb{R}^{K \times G}$ (softmax-normalized per topic) reconstructs gene probabilities as $P(\text{gene} \mid \theta) = \theta^\top \beta$, providing an explicit topic--gene mapping.

\textbf{Topic-FM-Transformer} replaces the MLP encoder with multi-head self-attention (4 heads) using a cell-as-token design, attending to heterogeneous cell properties within each batch. This captures cell--cell interaction patterns and yields the strongest NMI/ARI/ASW.

\textbf{Topic-FM-Contrastive} adds a MoCo-v2 contrastive head with data augmentation (Bernoulli dropout, noise injection) and a 4096-entry memory queue ($\tau{=}0.2$), combining instance-level contrast with the Dirichlet prior for the highest external win rate and best DAV.

\textbf{Topic-FM-GAT} employs a Graph Attention Network encoder operating over a precomputed $k$-nearest-neighbor graph ($k{=}15$). Two GATConv layers with 4 attention heads, residual connections, and LayerNorm aggregate neighbourhood information, enabling the model to exploit local transcriptomic similarity structure.

KL weight is fixed at 0.01 with free bits to prevent posterior collapse. Dirichlet sparsity strength $\kappa{=}10$. Pure-VAE, Pure-Transformer-VAE, and Pure-Contrastive-VAE serve as matched baselines.

\subsection{Flow Matching Refinement}
After the VAE warmup phase ($T_\text{warm}{=}50$ epochs), a conditional optimal-transport flow field is jointly trained in the pre-softmax space $\mathbb{R}^K$. Given a posterior sample $\log\theta$ (target) and Gaussian noise $\epsilon \sim \mathcal{N}(0,I)$ (source), the OT interpolant is $z_t = (1{-}t)\epsilon + t\log\theta$, with velocity target $v^* = \log\theta - \epsilon$. A small MLP velocity network $v_\psi(z_t, t)$ is trained to minimise
\[
  \mathcal{L}_\text{FM} = \mathbb{E}_{t\sim U(0,1)} \| v_\psi(z_t, t) - v^* \|^2.
\]
The total loss is $\mathcal{L} = \mathcal{L}_\text{ELBO} + \lambda_\text{FM}\,\mathcal{L}_\text{FM}$, with $\lambda_\text{FM}{=}0.1$. At inference, Euler integration from $t_0{=}0.8$ to $t{=}1.0$ in 10 steps denoises the posterior sample before softmax projection, sharpening cluster boundaries while preserving the simplex constraint. The flow field uses sinusoidal time embeddings (dim\,=\,32) and a 2-layer MLP ($128 \to 128$) with LayerNorm and Mish activations.

\subsection{Training and Data}
Canonical refresh settings are KL weight 0.01, dropout 0.0, weight decay $1\times10^{-3}$, learning rate $1\times10^{-3}$, 1000 epochs, and batch size 128. Flow matching activates at epoch 50 with weight $\lambda_\text{FM}{=}0.1$. Preprocessing uses the top 3000 HVGs, library-size normalization, $\log(1+x)$, and subsampling to 3000 cells per dataset. The refreshed internal benchmark spans 56 datasets (16 core plus 40 additional preprocessed cohorts).

\subsection{Evaluation}
We employ a multi-metric evaluation protocol centered on four core clustering metrics:

\textbf{Label concordance.}
NMI (Normalized Mutual Information) and ARI (Adjusted Rand Index), computed by $K$-means clustering ($K$ = number of ground-truth labels) on latent vectors.

\textbf{Geometric structure.}
ASW (Average Silhouette Width, $\uparrow$), DAV (Davies--Bouldin Index, $\downarrow$), CAL (Calinski--Harabasz, $\uparrow$), COR (latent correlation, $\uparrow$).

DRE (Dimensionality Reduction Evaluation) and LSE (Latent Structure Evaluation) families with their DREX/LSEX extensions are computed for reference but \emph{excluded from Topic-series composite scores}. The simplex-constrained latent space ($\Delta^{K-1}$) has bounded pairwise distances, rank-deficient covariance, and restricted manifold dimensionality, making distance-preservation metrics non-comparable to continuous-space reconstructions.

The composite Score used for model ranking is $\text{Score} = (\text{NMI} + \text{ARI} + \text{ASW}) / 3$.

\section{Results}\label{sec:results}

\subsection{Architecture Overview}
Figure~\ref{fig:arch} shows the four Topic-FM architectures with the shared flow matching refinement module.

\begin{figure}[H]
\centering
\includegraphics[width=\textwidth]{Fig1_arch_topic}
\caption{\textbf{Topic-FM architecture overview.} Four architectural variants share a common Dirichlet-prior VAE backbone ($n_\text{topics}$ = 10) with flow matching refinement in pre-softmax $\mathbb{R}^K$ space. (a)~Topic-FM-Base with MLP encoder; (b)~Topic-FM-Transformer with self-attention encoder; (c)~Topic-FM-Contrastive with MoCo contrastive head; (d)~Topic-FM-GAT with graph attention encoder over $k$NN. Dashed pink box indicates the shared flow matching module.}
\label{fig:arch}
\end{figure}

\subsection{Full-Catalogue Refresh}
Table~\ref{tab:topic_refresh} summarizes the mean performance of the Topic-FM models across 56 datasets.

\begin{table}[H]
\centering
\caption{\textbf{56-dataset Topic-FM refresh summary.} Mean metric values across \texttt{crossdata/results\_combined\_full}. Higher is better except DAV. DRE-UMAP and LSE shown for reference only (excluded from composite due to simplex constraint). Composite = (NMI\,+\,ARI\,+\,ASW)\,/\,3.}
\label{tab:topic_refresh}
\begin{tabular}{lccccccc}
\toprule
Model & NMI & ARI & ASW & DAV$\downarrow$ & DRE-UMAP$^\dagger$ & LSE$^\dagger$ & Composite \\
\midrule
Pure-VAE & 0.555 & 0.340 & 0.408 & 0.803 & 0.776 & 0.548 & 0.434 \\
Topic-FM-Base & 0.582 & 0.397 & 0.453 & 0.860 & 0.733 & 0.493 & 0.477 \\
Topic-FM-Transformer & \textbf{0.601} & 0.410 & \textbf{0.496} & 0.810 & 0.730 & 0.492 & \textbf{0.502} \\
Topic-FM-Contrastive & 0.596 & \textbf{0.411} & 0.462 & \textbf{0.802} & 0.746 & 0.489 & 0.490 \\
\bottomrule
\multicolumn{8}{l}{\footnotesize $^\dagger$DRE/LSE excluded from composite; simplex constraint limits comparability (see Methods).}
\end{tabular}
\end{table}

The defining characteristic of the flow-matching-refined Dirichlet prior is \emph{balanced, concordant improvement}: Topic-FM-Transformer raises NMI by 8.2\% (0.555 $\to$ 0.601), ARI by 20.4\% (0.340 $\to$ 0.410), and ASW by 21.7\% (0.408 $\to$ 0.496) simultaneously relative to Pure-VAE, yielding a composite gain of 15.6\% (0.434 $\to$ 0.502). This is a win across all three core metrics. Topic-FM-Contrastive achieves the best DAV (0.802) and highest ARI (0.411), while Topic-FM-Transformer leads on NMI, ASW, and composite score.

\subsection{Cross-Dataset Patterns}
Figure~\ref{fig:crossdataset} presents per-dataset scatter plots showing the balanced improvement at dataset granularity.

\begin{figure}[H]
\centering
\includegraphics[width=\textwidth]{Fig5_crossdataset_topic}
\caption{\textbf{Cross-dataset metric relationships.} Per-dataset scatter plots for concordance--geometry pairs (NMI vs.\ ASW, NMI vs.\ DAV, ARI vs.\ ASW, ARI vs.\ DAV) across all internal models on 56 datasets. Topic-FM variants consistently occupy the high-concordance, high-geometry quadrant, confirming balanced improvement across the full dataset range.}
\label{fig:crossdataset}
\end{figure}

The scatter plots confirm that the balanced improvement is not driven by a few favorable datasets. Topic-FM variants consistently position in the upper-right quadrant (high concordance \emph{and} high geometry), while Pure-VAE points cluster at lower concordance. The concordance--geometry correlation is positive for Topic-FM models, confirming that the Dirichlet prior's improvements on one metric axis do not come at the expense of the other.

\subsection{Internal Ablation}
Figure~\ref{fig:ablation} presents the curated ablation analysis with organized panels focusing on the metrics where Topic-FM models demonstrate their advantage.

\begin{figure}[H]
\centering
\includegraphics[width=\textwidth]{ablation/figures/uniform_grid}
\caption{\textbf{Internal ablation: Topic-FM vs.\ prior-free VAEs.} All 34 metrics across 6 evaluation suites (Clustering, DRE-UMAP, DRE-tSNE, LSE, DREX, LSEX) displayed on a uniform grid. Boxplots show metric distributions across 16 core datasets for the Topic-FM models and matched baselines.}
\label{fig:ablation}
\end{figure}

The ablation confirms that Topic-FM models achieve the best medians on all four concordance--geometry metrics: Topic-FM-Transformer or Topic-FM-Contrastive lead on NMI, ARI, ASW, and DAV. The red-bordered boxes (best median) fall on Topic-FM variants for the majority of panels. Calinski--Harabasz (CAL) is also higher for Topic-FM models, indicating tighter between-cluster separation.

\subsection{Statistical Validation}
Table~\ref{tab:topic_wilcoxon_core} presents Wilcoxon signed-rank tests confirming the statistical significance of the balanced improvements.

\begin{table}[!htbp]
\centering
\caption{\textbf{Wilcoxon signed-rank tests on core metrics (Topic series).} Pairwise structured vs.\ pure comparisons across 56 datasets. W/L/T = wins/losses/ties; Sig.\ at $\alpha{=}0.05$.}
\label{tab:topic_wilcoxon_core}
\begin{tabular}{llccccc}
\toprule
Structured vs Pure & Metric & W/L/T & $p$-value & Sig. & Cliff's $\delta$ & Effect \\
\midrule
Topic-Base vs Pure-VAE & ARI & 10/2/0 & 0.001 & Yes & 0.458 & medium \\
Topic-Base vs Pure-VAE & ASW & 11/1/0 & 0.013 & Yes & 0.569 & large \\
Topic-Base vs Pure-VAE & DAV & 4/8/0 & 0.849 & No & -0.236 & small \\
Topic-Base vs Pure-VAE & NMI & 11/1/0 & 7.32e-04 & Yes & 0.667 & large \\
Topic-Contrastive vs Pure-Contrastive-VAE & ARI & 12/0/0 & 2.44e-04 & Yes & 0.806 & large \\
Topic-Contrastive vs Pure-Contrastive-VAE & ASW & 12/0/0 & 2.44e-04 & Yes & 1.000 & large \\
Topic-Contrastive vs Pure-Contrastive-VAE & DAV & 12/0/0 & 2.44e-04 & Yes & 1.000 & large \\
Topic-Contrastive vs Pure-Contrastive-VAE & NMI & 12/0/0 & 2.44e-04 & Yes & 0.833 & large \\
Topic-Transformer vs Pure-Transformer-VAE & ARI & 8/4/0 & 0.026 & Yes & 0.403 & medium \\
Topic-Transformer vs Pure-Transformer-VAE & ASW & 10/2/0 & 0.002 & Yes & 0.722 & large \\
Topic-Transformer vs Pure-Transformer-VAE & DAV & 7/5/0 & 0.259 & No & 0.222 & small \\
Topic-Transformer vs Pure-Transformer-VAE & NMI & 11/1/0 & 4.88e-04 & Yes & 0.611 & large \\
\bottomrule
\end{tabular}
\end{table}


The balanced improvement is statistically robust. Topic-FM-Contrastive vs.\ Pure-Contrastive-VAE achieves 12/0/0 wins on all four core metrics ($p = 2.44\times10^{-4}$, Cliff's $\delta = 0.81$--$1.00$, all large effects). Topic-FM-Base and Topic-FM-Transformer also reach significance on NMI, ARI, and ASW (medium-to-large effects). DAV improvements reach significance only for Topic-FM-Contrastive, consistent with its role as the geometry specialist. The key result is that concordance and geometry significance are achieved simultaneously---confirming the balanced nature of the improvement.

\subsection{Downstream Classification}
Table~\ref{tab:topic_knn} reports $k$NN classification performance, providing a practical downstream evaluation that directly tests the learned representations.

\begin{table}[!htbp]
\centering
\caption{\textbf{$k$NN downstream classification (Topic-FM series).} Mean accuracy and macro-F1 across 12 datasets. Best value per column in bold.}
\label{tab:topic_knn}
\begin{tabular}{lcc}
\toprule
Model & Accuracy & F1-macro \\
\midrule
Topic-FM-Transformer & \textbf{0.733} & \textbf{0.517} \\
Topic-FM-Base & 0.729 & 0.506 \\
Topic-FM-Contrastive & 0.728 & 0.488 \\
Pure-VAE & 0.646 & 0.405 \\
Pure-Transformer-VAE & 0.617 & 0.368 \\
Pure-Contrastive-VAE & 0.603 & 0.355 \\
\bottomrule
\end{tabular}
\end{table}


Topic-FM-Transformer achieves the best accuracy (0.733) and macro-F1 (0.517), outperforming all three prior-free baselines by a substantial margin (Pure-VAE: 0.646/0.405, Pure-Contrastive-VAE: 0.603/0.355). All Topic-FM variants outperform all Pure-VAE variants, indicating that the flow-matching-refined Dirichlet prior consistently improves the discriminative quality of learned representations for classification tasks.

\subsection{Sensitivity Analysis}
Figure~\ref{fig:sensitivity} shows hyperparameter sensitivity for Topic-FM-Base.

\begin{figure}[H]
\centering
\includegraphics[width=\textwidth]{Fig3_sensitivity_topic}
\caption{\textbf{Hyperparameter sensitivity (Topic-FM-Base).} Nine hyperparameters swept across 12 datasets and four core metrics (NMI, ARI, ASW, DAV). Performance remains stable across a wide range of hyperparameter values, confirming that the flow-matching-refined Dirichlet prior is robust to hyperparameter choices.}
\label{fig:sensitivity}
\end{figure}

\subsection{UMAP Sweep Evolution}
Figure~\ref{fig:umap-sweep} shows UMAP evolution across sweeps.

\begin{figure}[H]
\centering
\includegraphics[width=\textwidth]{Fig4_training_topic}
\caption{\textbf{UMAP evolution across hyperparameter sweeps (Topic-FM-Base).} KMeans-colored embedding grids showing how latent geometry evolves across sweep values.}
\label{fig:umap-sweep}
\end{figure}

\subsection{Biological Validation and Interpretability}

A key advantage of the Dirichlet prior is that the decoder $\beta$ matrix provides an explicit gene--topic mapping, enabling dual-pathway biological validation: (i) perturbation-based gene-importance scoring, which identifies genes whose ablation most affects each latent component, and (ii) decoder-$\beta$-based analysis, which reads gene-program memberships directly from the learned parameters. This dual-pathway validation is unique to Dirichlet-prior models; methods without a topic--word decoder matrix must rely solely on perturbation-based importance. The flow matching refinement preserves this interpretability: it operates in pre-softmax space and does not alter the decoder $\beta$ weights.

\begin{figure}[H]
\centering
\includegraphics[width=\textwidth]{Fig6_biological_topic}
\caption{\textbf{Biological validation: perturbation gene-importance and decoder $\boldsymbol{\beta}$.} Z-scored perturbation importance and topic--gene decoder weights across representative datasets (endo, dentate). Perturbation heatmaps show gene importance per component; decoder $\beta$ heatmaps show the direct gene-program mapping unique to Dirichlet-prior models.}
\label{fig:bio}
\end{figure}

\begin{figure}[H]
\centering
\includegraphics[width=\textwidth]{Fig7_correlation_topic}
\caption{\textbf{Latent--gene Pearson correlation heatmaps (Topic-FM).} Top 30 genes per topic by $|r|$ across three datasets. Block-diagonal structure indicates that each topic captures a distinct expression program.}
\label{fig:corr}
\end{figure}

\begin{figure}[H]
\centering
\includegraphics[width=\textwidth]{Fig8_umap_topic}
\caption{\textbf{Latent UMAP projections (Topic-FM).} Component UMAP projections with cell-type overlays and top gene-expression overlays.}
\label{fig:umap}
\end{figure}

\begin{figure}[H]
\centering
\includegraphics[width=\textwidth]{Fig9_enrichment_topic}
\caption{\textbf{GO Biological Process enrichment (Topic-FM).} Top 20 salient genes per topic, including both perturbation-based and decoder-$\boldsymbol{\beta}$-based enrichment panels. Convergent enrichment across both pathways provides dual evidence that topics represent genuine gene regulatory programs.}
\label{fig:enrichment}
\end{figure}

Figures~\ref{fig:bio}--\ref{fig:enrichment} confirm that the flow-matching-refined Dirichlet prior encourages biologically coherent topic structure. Perturbation importance (Fig.~\ref{fig:bio}) and decoder $\beta$ weights independently identify distinct gene sets per topic; latent--gene correlations (Fig.~\ref{fig:corr}) show that individual topics track specific expression programs (top 30 genes per topic by $|r|$). UMAP projections (Fig.~\ref{fig:umap}) demonstrate spatial separation of known cell types through topic-weighted coloring. GO enrichment (Fig.~\ref{fig:enrichment}) yields significant biological-process terms (Fisher's exact, Bonferroni $p < 0.05$) under both the perturbation and decoder-$\beta$ pathways, providing convergent evidence that the learned topics represent genuine gene regulatory programs.

\subsection{External Benchmark}
Figure~\ref{fig:external} presents the unified external comparison.

\begin{figure}[H]
\centering
\includegraphics[width=\textwidth]{vs_external/figures/proposed/uniform_grid}
\caption{\textbf{External benchmark (proposed models).} All 34 metrics on a uniform grid comparing the four Topic-FM internal variants across 16 core datasets. Each boxplot shows the distribution across datasets with significance brackets.}
\label{fig:external}
\end{figure}

Table~\ref{tab:topic_external_winrate} quantifies the external win rates.

\begin{table}[!htbp]
\centering
\caption{\textbf{External win-rate summary (Topic-FM series).} Each variant is compared against 23 external baselines via Wilcoxon signed-rank tests across 56 datasets. Core win rate denotes the fraction of significant wins on the core metrics. Best value per column is shown in bold.}
\label{tab:topic_external_winrate}
\setlength{\tabcolsep}{5pt}
\small
\resizebox{\textwidth}{!}{%
\begin{tabular}{lccccc}
\toprule
Model & Total tests & Sig.\ (all) & Core sig. & Core total & Core win rate \\
\midrule
Topic-FM-Contrastive & 407 & 173 & \textbf{38} & 44 & \textbf{0.864} \\
Topic-FM-Transformer & 407 & 161 & 37 & 44 & 0.841 \\
Topic-FM-Base & 407 & 163 & 36 & 44 & 0.818 \\
Pure-VAE & 407 & \textbf{217} & 32 & 44 & 0.727 \\
Pure-Transformer-VAE & 407 & 177 & 27 & 44 & 0.614 \\
Pure-Contrastive-VAE & 407 & 148 & 26 & 44 & 0.591 \\
\bottomrule
\end{tabular}
}
\end{table}


Topic-FM-Contrastive achieves the highest core win rate (86.4\%), winning 38 of 44 core-metric comparisons. All Topic-FM variants outperform all Pure-VAE variants on core win rate, confirming that the flow-matching-refined Dirichlet prior's advantage generalizes to external baselines. The per-metric breakdown reveals ASW as the strongest axis (consistent near-universal wins against external models), while NMI shows the fewest significant wins, with CLEAR and scVI achieving higher mean concordance. No single external model matches the combined concordance--geometry--interpretability--downstream balance of the Topic-FM series.

\subsection{Computational Efficiency}
Table~\ref{tab:topic_runtime} reports the computational cost.

\begin{table}[!htbp]
\centering
\caption{\textbf{Runtime analysis (Topic-FM series).} Mean training time, throughput, peak GPU memory, and parameter count across 56 datasets. Best value per column in bold.}
\label{tab:topic_runtime}
\begin{tabular}{lcccc}
\toprule
Model & Time (s) & Sec/Epoch & GPU (MB) & Params \\
\midrule
Pure-VAE & \textbf{45.0} & \textbf{0.0450} & \textbf{825.3} & \textbf{433,476} \\
Topic-FM-Base & 45.6 & 0.0456 & 903.8 & 433,476 \\
Topic-FM-Transformer & 94.1 & 0.0941 & 936.9 & 3,440,068 \\
Pure-Transformer-VAE & 101.3 & 0.1013 & 910.8 & 3,440,068 \\
Topic-FM-Contrastive & 111.7 & 0.1117 & 923.8 & 435,784 \\
Pure-Contrastive-VAE & 113.6 & 0.1136 & 873.8 & 435,784 \\
\bottomrule
\end{tabular}
\end{table}


Topic-FM-Base adds minimal overhead over Pure-VAE (the Dirichlet prior's KL computation and flow matching velocity field training are both efficient). Topic-FM-Transformer is $\sim$2$\times$ slower due to self-attention. The flow matching module adds a small 2-layer MLP ($128 \to 128$, sinusoidal time embedding dim\,=\,32) that trains only after warmup and requires only 10 Euler steps at inference.

\section{Discussion}\label{sec:discussion}

\textbf{Balanced improvement without trade-off.}
The defining result of the Topic-FM series is that the flow-matching-refined Dirichlet prior improves label concordance \emph{and} cluster geometry simultaneously. Topic-FM-Transformer raises all three core metrics over Pure-VAE: NMI +8.2\%, ARI +20.4\%, ASW +21.7\%, yielding a composite gain of 15.6\%. Wilcoxon tests confirm significance with medium-to-large effects on all three metrics. The balanced profile makes Topic-FM the practical default when both clustering quality and structural coherence are valued.

\textbf{Flow matching as latent refinement.}
The conditional OT flow field operates in pre-softmax $\mathbb{R}^K$ space, denoising posterior samples via partial Euler integration ($t_0{=}0.8 \to 1.0$). This refinement sharpens cluster boundaries in the latent space while preserving the simplex constraint after softmax projection. Crucially, the flow field is trained with a simple MSE velocity loss and activates only after the VAE has converged, avoiding interference with early-stage posterior optimisation.

\textbf{Variant specialization.}
Table~\ref{tab:topic_variant_ranking} provides the full individual variant ranking.
Topic-FM-Transformer achieves the highest composite score (0.502) and leads on NMI (0.601) and ASW (0.496), benefiting from self-attention's ability to capture cell--cell correlations. Topic-FM-Contrastive achieves the highest ARI (0.411), best DAV (0.802), and the highest external win rate (86.4\%), confirming that contrastive learning synergizes with the Dirichlet prior for the strongest generalization. Topic-FM-Base provides the simplest architecture with a strong 9.8\% composite improvement and negligible computational overhead. Topic-FM-GAT leverages local graph structure via $k$NN-based attention, offering an alternative to batch-level self-attention.

\begin{table}[!htbp]
\centering
\caption{\textbf{Individual variant ranking (Topic-FM series).} Models ranked by composite score across 56 datasets. Best value per column in bold.}
\label{tab:topic_variant_ranking}
\begin{tabular}{lccccc}
\toprule
Model & NMI & ARI & ASW & DAV $\downarrow$ & Score \\
\midrule
Topic-FM-Transformer & \textbf{0.569} & \textbf{0.342} & \textbf{0.458} & 0.828 & \textbf{0.457} \\
Topic-FM-Contrastive & 0.560 & 0.335 & 0.418 & \textbf{0.801} & 0.438 \\
Topic-FM-Base & 0.562 & 0.324 & 0.416 & 0.897 & 0.434 \\
Pure-VAE & 0.456 & 0.230 & 0.359 & 0.855 & 0.348 \\
Pure-Transformer-VAE & 0.457 & 0.252 & 0.321 & 0.921 & 0.343 \\
Pure-Contrastive-VAE & 0.397 & 0.160 & 0.231 & 1.184 & 0.263 \\
\bottomrule
\end{tabular}
\end{table}


\textbf{Downstream classification confirms practical utility.}
The $k$NN classification results provide independent evidence that the flow-matching-refined Dirichlet prior improves representation quality for practical tasks. Topic-FM-Transformer's accuracy (0.733) and F1 (0.517) exceed all Pure-VAE variants by 13--22\%, demonstrating that the balanced metric improvement extends to discriminative downstream use.

\textbf{Interpretability as a built-in feature.}
The Dirichlet VAE contains an explicit topic--word matrix $\beta$ that maps each topic to gene-expression probabilities. This enables dual-pathway biological validation: perturbation-based importance identifies genes that affect each topic, while $\beta$ weights directly read out gene-program membership. Both pathways yield significant GO enrichment (Bonferroni $p < 0.05$) on endoderm and dentate datasets, confirming that the learned topics are biologically meaningful. The flow matching module preserves this interpretability: it refines the pre-softmax representation without modifying the decoder $\beta$ weights.

\textbf{Cross-dataset robustness.}
The scatter-plot analysis (Fig.~\ref{fig:crossdataset}) demonstrates that Topic-FM models' balanced improvement is not an artifact of a few favorable datasets. The concordance--geometry scatter shows Topic-FM variants in the upper-right quadrant (high on both axes) across all 56 datasets, with a positive concordance--geometry correlation that confirms the Dirichlet prior avoids the trade-off pattern.

\textbf{External positioning.}
Among external baselines, scVI is the sole model with higher NMI/ARI than Topic-FM-Transformer, and scDAC is the sole model with higher ASW. Neither matches the combined concordance--geometry--interpretability--downstream profile of Topic-FM models. Topic-FM-Contrastive's 86.4\% core win rate establishes it as the strongest overall model in the external comparison.

\textbf{Limitations.}
The simplex constraint makes DRE/LSE metrics non-comparable across prior families and limits the latent space to $K{-}1$ effective dimensions. The number of topics ($K{=}10$) is a fixed hyperparameter that may not suit all datasets. The flow matching warmup period ($T_\text{warm}{=}50$) adds a hyperparameter, though sensitivity analysis shows robustness across a range of values. External baselines with zero performance on many datasets (6 of 23 models) were excluded from win-rate calculations.

\section*{Declarations}
Funding, competing interests, and data/code availability to be completed by authors.

\reftitle{References}
\bibliography{sn-bibliography}

\end{document}
